\documentclass[10.5pt]{article}
\usepackage{amsmath, amsfonts, amssymb,amsthm}
\usepackage[includeheadfoot]{geometry} % For page dimensions
\usepackage{fancyhdr}
\usepackage{enumerate} % For custom lists
\usepackage{tikz-cd}
\usepackage{graphicx}

\fancyhf{}
\lhead{MAT1002 hw4}
\rhead{Tighe McAsey - 1008309420}
\pagestyle{fancy}

% Page dimensions
\geometry{a4paper, margin=1in}

\theoremstyle{definition}
\newtheorem{pb}{}
\usepackage{tikz-cd, stackengine}

% Commands:

\newcommand{\set}[1]{\{#1\}}
\newcommand{\gen}[1]{\langle#1\rangle}
\newcommand{\abs}[1]{\left\vert#1\right\vert}
\newcommand{\norm}[1]{\lvert\lvert#1\rvert\rvert}
\newcommand{\tand}{\text{ and }}
\newcommand{\tor}{\text{ or }}
\newcommand{\pd}{\frac{\partial}{\partial x_j}}
\newcommand{\px}{\frac{\partial}{\partial x}}
\newcommand{\py}{\frac{\partial}{\partial y}}
\newcommand{\pz}{\frac{\partial}{\partial z}}
\newcommand{\pzb}{\frac{\partial}{\partial \bar{z}}}
\newcommand{\ppx}{\frac{\partial^2}{\partial x^2}}
\newcommand{\ppy}{\frac{\partial^2}{\partial y^2}}
\newcommand{\ppz}{\frac{\partial^2}{\partial z^2}}
\newcommand{\hess}{\operatorname{Hess}}
\newcommand{\Log}{\operatorname{Log}}
\newcommand{\Arg}{\operatorname{Arg}}
\newcommand{\dv}{dz\wedge d\bar{z}}
\newcommand{\pzpz}{\frac{\partial^2}{\partial z \partial \bar{z}}}
\newcommand{\pxpx}{\frac{\partial^2}{\partial x^2}}
\newcommand{\pypy}{\frac{\partial^2}{\partial y^2}}
\newcommand{\pwpw}{\frac{\partial^2}{\partial w_0 \partial \bar{w_0}}}

\begin{document}
    \begin{pb}
        \textbf{(i)} Suppose \(T^\dagger v = 0\), then for any \(w \in V_1\), we have \(0 = \gen{T^\dagger v, w} = \gen{v, Tw}\), so \(v \in (\operatorname{Im} T)^\perp\).  Conversely if \(v \in (\operatorname{Im} T)^\perp\), then \(\gen{T^\dagger v, -} \equiv 0\), and since \(\gen{-,-}\) is an inner product it is non-degenerate implying that \(T^\dagger v = 0\). The proof for \(S\) is identical.

        \textbf{(ii)} Suppose \(u = T^\dagger v\), then for any \(w \in \ker T\) we find that \(\gen{u,w} = \gen{T^\dagger v,w} = \gen{v,Tw} = 0\) hence \(u \in (\ker T)^\perp\), this implies that \(\operatorname{Im}T^\dagger \subset (\ker T)^{\perp}\). Then to see they are equal it suffices to show their dimensions are equal, we first use rank nullity so that \(\dim (\ker T)^\perp = \dim V_1 - \dim \ker T = \operatorname{rank} T\), so we only need show that \(\operatorname{rank} T = \operatorname{rank} T^\perp\). To show this we write
        \begin{align}
            \dim V_2 &= \dim(\ker T^\dagger) + \operatorname{rank}T^\dagger \\
            \dim V_2 &= \dim(\operatorname{Im} T)^\perp + \operatorname{rank} T
        \end{align}
        This gives us
        \begin{align}
            \dim(\ker T^\dagger) + \operatorname{rank}T^\dagger = \dim(\operatorname{Im} T)^\perp + \operatorname{rank} T
        \end{align}
        then applying part (i), we have \(\ker T^\dagger = (\operatorname{Im} T)^\perp\), so cancelling these terms we get \(\operatorname{rank} T^\dagger = \operatorname{rank} T\) allowing us to conclude that indeed \(\operatorname{Im}T^\dagger\), impling that \(\dim \operatorname{Im} T^\dagger = \dim (\ker T)^\perp\) so that the vectorspaces must be equal. The same argument applies to \(S\).

        \textbf{(iii)} Let \(v \in \ker S \cap (\operatorname{Im} T)^\perp\), then for any \(w \in V^2\), we have
        \begin{align}
            \gen{\Delta v, w} &= \gen{TT^{\dagger}v,w} + \gen{S^\dagger S v, w} \overset{\ker S}{=} \gen{TT^{\dagger}v,w}\\ 
            &= \overline{\gen{w,TT^{\dagger}v}} = \overline{\gen{T^\dagger w, T^\dagger v}} = \gen{T^\dagger v, T^\dagger w} = \gen{v,TT^\dagger w} \overset{(\operatorname{Im} T)^\perp}{=} 0
        \end{align}
        so that \(\gen{\Delta v,-} \equiv 0\) which implies \(\Delta v = 0\) by non-degeneracy.

        Conversely, suppose that \(\Delta v = 0\), then
        \begin{align}
            0 = \norm{\Delta v}^2 = \norm{TT^\dagger v}^2 + \norm{S^\dagger S v}^2 + \gen{S^\dagger S v, TT^\dagger v} + \overline{\gen{S^\dagger S v, TT^\dagger v}}
        \end{align}
        but then \(\gen{S^\dagger S v, TT^\dagger v} = \gen{Sv,STT^{\dagger}v} = 0\) since \(ST = 0\), so that \(\norm{TT^\dagger v}^2 = 0 \tand \norm{S^\dagger S v}^2 = 0\). Now since \(\ker S^\dagger = (\operatorname{Im} S)^\perp\), and \(S^\dagger S v = 0\), we have \(Sv \in (\operatorname{Im} S)^\perp \cap \operatorname{Im} S = 0\), so that \(v \in \ker S\). Now since \(TT^\dagger v = 0\), we have 
        \begin{align}
            T^\dagger v \in \ker T \cap \operatorname{Im} T^\dagger \overset{(ii)}{=} \ker T \cap (\ker T)^\perp = 0
        \end{align}
        so that \(v \in \ker T^\dagger \overset{(i)}{=} (\operatorname{Im} T)^\perp\). Since \(v\) lies in both \((\operatorname{Im} T)^\perp\) and \(\ker S\), it lies in the intersection as desired.

        \textbf{(iv)} We first note that
        \begin{align}
            \ker \Delta = \operatorname{Im} T^{\perp} \cap \ker S = \operatorname{Im} T^{\perp} \cap (\operatorname{Im} S^\dagger)^\perp = (\operatorname{Im} T + \operatorname{Im} S^\dagger)^\perp
        \end{align}
        so that the decomposition \(V_2 = \ker \delta \oplus (\operatorname{Im} T + \operatorname{Im} S^\dagger)\) is orthogonal, it remains to check that \(\operatorname{Im} T\) and \(\operatorname{Im} S^\dagger\) are orthogonal, this is simple since for any \(v \in V_1\) and \(w \in V_3\) we have \(\gen{S^\dagger w, Tv} = \gen{w,ST v} = 0\) since \(ST = 0\). Thus the decomposition is orthogonal.

        Now we can define \(\ker \Delta \to \ker S/\operatorname{Im} T\) as \(v \mapsto [v]\), since \(\ker \Delta \subset \ker S\), to check this is an isomorphism assume that \(v \in \operatorname{ker} \Delta \cap \operatorname{Im} T\), then \(v = 0\), since \(\ker \Delta \subset (\operatorname{Im} T)^\perp\), so that this map is injective. Now, let \(w \in \ker S\), we want to show that \([w]\) is in the image of our map, so it suffices to show that for some \(u \in \operatorname{Im} T\) we have \(w - u \in \operatorname{Im} \Delta\). We take \(u\) to be the projection of \(w\) onto \(\operatorname{Im} T\), then \(w - u \in \ker S \cap (\operatorname{Im} T)^\perp = \ker \Delta\), so the map is onto and \(\ker \Delta \cong \ker S/ \operatorname{Im} T\) \qed
    \end{pb}
    \begin{pb}
        \textbf{(i)} We calculate the laplacian on \(\phi_{m,n}\)
        \begin{align}
            \Delta \phi_{m,n} &= -\pxpx e^{2\pi i (mx + ny)} -\pypy e^{2\pi i (mx + ny)} = -(m^2 + n^2)e^{2\pi i (mx + ny)} = 4\pi^2(m^2 + n^2)\phi_{m,n}
        \end{align}

        \textbf{(ii)}
        \begin{align}
            \gen{\phi_{m,n},\phi_{p,q}} &= \int_{T^2}e^{2\pi i (mx + ny - px - qy)}dxdy = \int_0^1 e^{2\pi i(m - p)x}\int_0^1 e^{2\pi i(n - q)y}dydx \\
            &= \int_0^1 \delta_{n,q} e^{2\pi i(m - p)x} dx = \delta_{n,q}\delta_{m,p}
        \end{align}

        \textbf{(iii)} Enumerate \(\mathbb{Z}^2\) as the set \(\set{r_1,r_2,r_3,\hdots}\), then we have \(f_n = \sum_1^n a_{r_j} \phi_{r_j} \) with \(\lim_{n\to\infty} f_n = f\), it follows that
        \begin{align}
            \lim_{m\to\infty}\gen{f_m,f_m} = \lim_{m\to\infty}\lim_{n\to\infty}\gen{f_m,f_n} = \gen{f_m,f} = \gen{f,f}
        \end{align}
        Where the first equality comes from the identity in (ii). It suffices to just compute \(\gen{f_m,f_m}\), but 
        \begin{align}
            \gen{f_m,f_m} = \sum_{1\leq i,j \leq m} \gen{a_{r_i}\phi_{r_i},a_{r_j}\phi_{r_j}} = \sum_{1\leq i,j \leq m} a_{r_i}\overline{a_{r_j}}\delta_{i,j} = \sum_1^n \abs{a_{r_i}}^2
        \end{align}
        from this we deduce that
        \begin{align}
            \gen{f,f} = \sum_{m,n\in \mathbb{Z}^2}\abs{a_{m,n}}^2
        \end{align}

        \textbf{(iv)} Since \(f\) is smooth, its partial derivatives are equal to the partials of its fourier series, whence we find that
        \begin{align}
            \int_{T^2}\abs{\nabla f}^2dxdy = \int_T \abs{\px f}^2 + \abs{\py f}^2 = \norm{\px f}^2 + \norm{\py f}^2
        \end{align}
        and the fourier coefficients of \(\px f\) corresponding to \(e^{2\pi i(mx + ny)}\) are \(2\pi i m a_{m,n}\), and for \(\py f\) we have \(2\pi i n a_{m,n}\), it follows that
        \begin{align}
            \int_{T^2}\abs{\nabla f}^2dxdy = 4\pi^2\left(\sum_{\mathbb{Z}^2} m^2\abs{a_{m,n}^2}  + \sum_{\mathbb{Z}^2}n^2\abs{a^2_{m,n}}\right) = 4\pi^2 \sum_{\mathbb{Z}^2} (m^2 + n^2)\abs{a_{m,n}}^2
        \end{align}

        Similarly, we can compute
        \begin{align}
            \abs{D^2 f}^2 = \abs{\pxpx f}^2 + \abs{\pypy f}^2 + \abs{\px\py f}^2 + \abs{\py\px f}^2
        \end{align}
        then we can compute the fourier coefficients of each of these partials, with the coefficient for \(\phi_{m,n}\) being \(-4\pi^2 m^2 a_{n,m}\) for \(\pxpx f\), \(-4\pi^2 n^2 a_{n,m}\) for \(\pypy f\) and \(-4\pi^2 nm a_{n,m}\)
        \begin{align}
            \int_{T^2}\abs{D^2 f}^2 &= \norm{\pxpx f}^2 + \norm{\pypy f}^2 + \norm{\px\py f}^2 + \norm{\py\px f}^2 \\
            &= 16\pi^4 \sum_{\mathbb{Z}^2} (n^4 + m^4 + 2n^2m^2)\abs{a_{n,m}}^2 = 16\pi^4 \sum_{\mathbb{Z}^2} (n^2 + m^2)^2 \abs{a_{n,m}}^2
        \end{align}

        \textbf{(v)} We can compute the fourier coefficients of \(f\) using that \(u \in W^{2,2}(T^2)\) has second weak derivatives. Then
        \begin{align}
            \int_{T^2} (\Delta u) e^{-2\pi i (nx + my)}dxdy &\overset{\text{IBP}}{=} -\int_{T^2}u \left(\pxpx + \pypy\right)e^{-2\pi i (nx + my)}dxdy \label{Laplacian1} \\
            &= \int_{T^2} u 4\pi^2 (n^2 + m^2)e^{-2\pi i (nx + my)}dxdy \\
            &= 4\pi^2(n^2 + m^2) a_{n,m} \label{Laplacian3}
        \end{align}
        So that \(f\) has fourier coefficients given by \(4\pi^2(n^2 + m^2) a_{n,m}\), it follows that \(\norm{f}_{L^2}^2 = 16\pi^4 \sum_{\mathbb{Z}^2}(n^2 + m^2)^2\abs{a_{n,m}}^2\), then
        \begin{align}
            \norm{u}_{w^{2,2}}^2 = \sum_{\mathbb{Z}^2} ((m^2 + n^2)^2 + 2(m^2 + n^2))\abs{a_{n,m}}^2 \leq \frac{1}{16\pi^4}\norm{f}^2_{L^2} + 2\norm{u}_{w^{1,2}}^2 \leq 2 \norm{f}^2_{L^2} + 2\norm{u}_{w^{1,2}}^2
        \end{align}
        so we can choose \(C = 2\).

        \textbf{(vi)} Let \(a_{n,m}\) be the Fourier coefficients of \(f\), then \(a_{0,0} = \int_{T^2} fdxdy = 0\), we take 
        \begin{align}
            u(x,y) := \sum_{\mathbb{Z}^2 \setminus \set{(0,0)}}\frac{1}{4\pi^2(n^2 + m^2)} a_{n,m} \phi_{n,m}(x,y) \label{definvlap}
        \end{align}
        Then the fourier coefficients of \(\Delta u\) are the fourier coefficients of \(f\) using \eqref{Laplacian1}-\eqref{Laplacian3}, since \(f \in L^2(T^2)\) it is equal to its Fourier series so that this \(u\) solves the Laplace equation. Moreover, we can check that \(u \in W^{2,2}\) directly
        \begin{align*}
            \norm{u}_{W^{2,2}}^2 = \sum_{\mathbb{Z}^2 \setminus \set{(0,0)}} \frac{(1 + m^2 + n^2)^2}{16\pi^4(n^2 + m^2)^2} \abs{a_{n,m}}^2 \leq \sum_{\mathbb{Z}^2} \abs{a_{n,m}}^2 = \norm{f}_{L^2}^2 < \infty
        \end{align*}
        Moreover since any \(u\) satisfying \(\Delta u = f\) must have the same Fourier coefficients for \((n,m) \neq (0,0)\) by \eqref{Laplacian1}-\eqref{Laplacian3}, hence they have the same fourier series apart from the coefficient of \(e^0 = 1\), any choice of constant still satisfies \(\Delta u = f\), so the coset \(u + \mathbb{C}\) are all the possible solutions. We first check why any solution of the following form must take \(\gen{u,1} = 0\), computing explicitly we get
        \begin{align}
            \int_{T^2} u &= \int_{T^2}\int_{T^2} G(x_1-x_2,y_1-y_2)f(x_2,y_2)dx_2dy_2dx_1dy_1 \\ 
            &\overset{\text{Fubini}}{=} \int_{T^2}\int_{T^2} G(x_1-x_2,y_1-y_2)f(x_2,y_2)dx_1dy_1dx_2dy_2 \\
            &= \int_{T^2}f(x_2,y_2)\int_{T^2}G(x_1-x_2,y_1-y_2)dx_1dy_1dx_2dy_2 \\
            &\overset{\text{change of variables}}{=} \int_{T^2}f(x_2,y_2)\int_{T^2}G(x_1,y_1)dx_1dy_1dx_2dy_2 \\
            &= \left(\int_{T^2} f\right)\left(\int_{T^2} G\right) = 0
        \end{align}
        We can apply fubini since if \(G \in L^2\), then \(\norm{fG}_1 \leq \norm{f}_2 \norm{G}_2\) by Holders. Now since such a solution to \(\Delta u\) is unique since the constant is fixed, it suffices to find some \(G\), such that
        \begin{align}
            \Delta_{x_1,y_1} \int_{T^2} G(x_1-x_2,y_1-y_2)f(x_2,y_2)dx_2dy_2 = f(x_1,y_1) \label{desiderata}
        \end{align}
        We define \(G\) as follows, it is immediate that \(G \in L^2\), and by Holder's inequality also \(L^1\) since \(T^2\) has finite measure.
        \begin{align}
            G(x,y) = \sum_{\mathbb{Z}^2 \setminus \set{(0,0)}} \frac{1}{4\pi ^2(m^2 + n^2)}\phi_{n,m}
        \end{align}
        Now we finally compute that \(G\) satisfies \eqref{desiderata}, to do so we compute the weak derivative \(\Delta G\), since finite linear combinations of \(\phi_{n,m}\) are dense in \(C^\infty(T^2)\) we know that distributions are uniquely determined by their ``Fourier coefficients'' so we can check that \(\Delta G = \delta_{0,0}\) the dirac delta in the weak sense. We can compute this as follows, (once again enumerating \(\mathbb{Z}^{2} \setminus \set{(0,0)}\)) we take \(S_N\) to be the \(N\)-th partial sum of Fourier coefficients, then \(S_N \to G\) in \(L^2\), so that by definition of weak derivative, \(\Delta G(1) = \gen{G,\Delta 1} = \gen{G,0} = 0\), and for \((n,m) \neq (0,0)\)
        \begin{align}
            \Delta G(\phi_{n,m}) = \gen{G, \Delta \phi_{n,m}} = 4\pi^2(m^2 + n^2)\gen{\lim_{N\to\infty}S_N,\phi_{n,m}} = \lim_{N\to\infty} 4\pi^2(m^2 + n^2) \gen{S_N,\phi_{n,m}} = 1
        \end{align}
        So that \(\Delta G = \delta_{(0,0)} - 1\) as a weak derivative since they act the same way on \(\phi_{n,m}\) which are dense in \(C^\infty(T^2)\), we conclude that
        \begin{align}
            \Delta_{x_1,y_1} \int_{T^2} G(x_1-x_2,y_1-y_2)f(x_2,y_2)dx_2dy_2 = (\delta_{x_1,y_1} - 1)(f) = f(x_1,y_1)
        \end{align} \qed
    \end{pb}
    \begin{pb}
        Since Riemann surfaces are compact, any global holomorphic function must attain its maximum, thus by the open mapping principle must be constant in the chart it attains its maximum, since this chart is open the global holomorphic function must be constant by the identity principle. Now considering the divisor \(0\), the previous discussion tells us that \(H^0(X,\mathcal{O}_X(0)) = \mathbb{C}\) has complex dimension 1. We use Riemann Roch:
        \begin{align}
            \dim_\mathbb{C}H^0(X,\mathcal{O}_X(D)) - \dim_\mathbb{C}H^0(X,K_X(-D)) = \deg D + 1 - g
        \end{align}
        where \(g\) is the genus, now since \(D = 0\) has degree zero, and \(\dim_\mathbb{C}H^0(X,\mathcal{O}_X(D)) = 1\) the equation simplifies to
        \begin{align}
            \dim_\mathbb{C}H^0(X,K_X(0)) = g
        \end{align}
        Where by definition \(H^0(X,K_X(0))\) is the space of global holomorphic differentials. Rewriting \(k = 2j + 1\) for \(j \in \mathbb{Z}_{\geq 0}\) we can recall from the cut and paste construction that the genus of \(X\) is \(g = j+1\), this also gives us \(g = \frac{k+1}{2}\). This suffices to show the following holomorphic differentials form a basis for \(w = \sqrt{z(z-1)(z-\lambda_1)\cdots(z-\lambda_{k})}\)
        \begin{align}
            \frac{dz}{w}, \frac{zdz}{w}, \hdots, \frac{z^{g-1}dz}{{w}}
        \end{align}
        These are trivially holomorphic away from \(\set{0,1,\lambda_i,\infty}\), to check at these points we can just plug in local coordinates at \(z = 0\) (since \(z = 1\) and \(\lambda_i\) behave the same) and local coordinates at \(\infty\). Near zero we get coordinates \(t\) with \(z = t^2\), with the signs of \(t\) consistent with the gluing of \(X\), in which case we get
        \begin{align}
            z^i dz/w = t^{2i} \frac{2tdt}{t\sqrt{(t^2 - 1)\cdots(t^2 - \lambda_k)}} = t^{2i} \frac{dt}{\sqrt{(t^2 - 1)\cdots(t^2 - \lambda_k)}}
        \end{align}
        is holomorphic. And near \(t = \infty\) we have coordinates given by \(t^2 = \frac{1}{z}\) once again with the sign of \(t\) gluing across \(X\), so that we get
        \begin{align}
            z^i dz/w = t^{-2i}\frac{-2dt}{t^3} \frac{t^{k+2}}{\sqrt{(1-t^2)(1-\lambda_1t^2)\cdots(1-\lambda_k t^2)}} = t^{2g + 1 - 2i - 3} dt \cdot (\text{holomorphic})
        \end{align}
        Then \(i \leq g-1\) gives us that \(2i + 3 \leq 2g+1\) so this is indeed a holomorphic differential. \qed
    \end{pb}
\end{document}